% Finalizado 0 
\begin{center}
    \textbf{Resumen}
\end{center}

\noindent El presente Trabajo de Fin de Grado describe el desarrollo de un motor gráfico implementado con Vulkan. El objetivo principal consiste en alcanzar un alto rendimiento en la simulación y el renderizado de partículas en tiempo real. Para tal fin se ha diseñado una arquitectura paralela basada en \textit{Compute Shaders} mediante la cual se emulan efectos volumétricos complejos como fuego, humo y agua. La piedra angular del sistema reside en la delegación íntegra del ciclo de vida de las partículas a la GPU. El procesador principal queda exento de actualizar posiciones, propiedades fotométricas y dinámicas físicas. Por consiguiente, se mitiga el tradicional cuello de botella en la transferencia de datos entre la CPU y la memoria de vídeo, limitación recurrente en interfaces gráficas de mayor nivel de abstracción.

El motor gráfico se fundamenta en un control exhaustivo sobre el hardware, estructurándose mediante módulos independientes. El \textit{pipeline} integra computación procedimental, sincronización explícita a través de barreras de memoria (\texttt{VkBufferMemoryBarrier}) y renderizado volumétrico mediante el uso conjunto de \textit{Point Sprites} y \textit{Alpha Blending}. Adicionalmente a los sistemas de partículas, la aplicación provee un entorno tridimensional equipado con un modelo de iluminación empírico ADS (Ambiente, Difusa y Especular), un intérprete optimizado de geometría y un \textit{Skybox} dinámico. El empleo de descriptores y búferes locales garantiza un acceso de baja latencia a los recursos. Los resultados experimentales confirman la eficiencia de Vulkan: el sistema es capaz de computar simultáneamente una gran cantidad de partículas, sosteniendo una tasa estable de 60 fotogramas por segundo en procesadores gráficos de gama media, al tiempo que ofrece amplias opciones de parametrización visual. En conclusión, el proyecto constituye una base arquitectónica sólida y escalable para la integración de simulaciones físicas en aplicaciones gráficas contemporáneas.                                      

\vspace{0.5cm}
\textit{Palabras clave:} Vulkan, compute shaders, sistemas de partículas, tiempo real, GPU

\vspace{0.5cm}
\textit{Competencias específicas:} 
En este trabajo se han desarrollado las siguientes competencias específicas:
\begin{itemize}
    \item Capacidad para diseñar, programar y evaluar sistemas interactivos que muestren información compleja. Aplicación práctica de estos conocimientos para resolver problemas reales dentro del diseño de interacción persona-computadora.
\end{itemize}

